\documentclass[10pt,a4paper]{amsart}
\usepackage[margin=19mm]{geometry}
\usepackage{amsmath,amssymb,mathtools}
\usepackage[scr=rsfs,cal=euler]{mathalfa}
\usepackage{xfrac}
\usepackage[unicode,hidelinks,bookmarks=false]{hyperref}
\newcommand{\ie}{i.e.\ }
\newcommand{\cf}{cf.\ }
\newcommand{\resp}{respectively\ }
\newcommand{\Subsection}{\subsection}
\providecommand{\nobreakdash}{\nobreak\hbox{-}}
\setlength{\emergencystretch}{2em}
\multlinegap0pt
\usepackage{natbib}
\usepackage{float}
\usepackage{amssymb}
\usepackage{caption}
\newtheorem{theorem}{Theorem}
\newcommand{\R}{\mbox{$\mathbb{R}$}}
\newcommand{\abs}[1]{\left\vert#1\right\vert}
\title{On Koebe-type functions for harmonic quasiconformal mappings}
\author{Zhi-Gang Wang, Jia-Le Qiu, Antti Rasila}
\date{Source: arXiv:2505.12562v3 (CC BY 4.0)}
\begin{document}
\maketitle

\noindent\emph{Source and licence.} On Koebe-type functions for harmonic quasiconformal mappings. Version arXiv:2505.12562v3 (CC BY 4.0), distributed by arXiv under the Creative Commons Attribution 4.0 International licence, http://creativecommons.org/licenses/by/4.0/, as recorded in the arXiv licence metadata for this exact version and shown on the abstract page https://arxiv.org/abs/2505.12562v3. Full licence notice: \texttt{licenses/arxiv-2505.12562-CC-BY-4.0.txt}.\par\smallskip

\noindent\textit{Editorial scope.} Counted unit: the article's theorem T4 (source0.tex lines 523--533). The article's complete original proof of it is delivered with it (source0.tex lines 534--588, one \texttt{proof} environment of 55 lines). No mathematics is written, repaired or re-derived: the statement and the proof are the article's own lines, in the article's own notation, and no retained line is substituted or re-flowed.  Retained uncounted as the source-local context the proof needs: lines 303--308; lines 453--454; lines 492--494; lines 495--497.  Nothing was omitted from any retained span, so the package carries no omission record.  No retained line contains a citation, so the wrapper carries no bibliography and no bibliography entry.  No retained line contains a figure environment, an image-inclusion command or drawing code, so no image is delivered and the package declares no asset dependency.  Editorial formatting only, in addition: an AMS \texttt{amsart} preamble that restates the article's own declarations and macros used by the retained lines, namely the environments \texttt{theorem} on separate counters, together with the standard packages those lines need.  The retained environments print in this document as Theorem 1 in order of appearance, whereas the article prints the same items with the numbering of its own full document.  The complete native source of the article as distributed in the arXiv e-print for this exact version is bundled unchanged as \texttt{sources/178322/source0.tex}.
\begin{equation}\label{1}
  \begin{cases}
      h(z)-g(z)=k_a(z)\quad(a\in\R),\\
     \,{g'(z)}/{h'(z)}=\lambda z\quad(0\leq\lambda<1),
  \end{cases}  
\end{equation}

\begin{equation}\label{T2-3} 	  k_a'(z)=\dfrac{{(1+z)}^{a-1}}{{(1-z)}^{a+1}}.
	\end{equation}

        \begin{equation}\label{T3-3}
            k_a''(z)=\dfrac{2(a+z)}{(1-z)^4}\left(\dfrac{1+z}{1-z}\right)^{a-2},
        \end{equation}

        \begin{equation}\label{T3-4}
            k_a'''(z)=\dfrac{2\left(1+z\right)^{a-3}}{\left(1-z\right)^{a+3}}\left(3z^2+6az+2a^2+1\right)
        \end{equation}

\begin{theorem}\label{T4}
	If $f=h+\bar{g}\in\mathcal{K}_{a,\lambda}$, then %$a_1=1$, 
    $$|a_2|\leq|a|+\frac{\lambda}{2},$$ 
	$$|a_3|\leq\dfrac{1}{3}\left(\lambda^2+2\abs{a}\lambda+2a^2+1\right),$$
	$$|a_4|\leq\dfrac{1}{3}\abs{a}^3+\dfrac{2}{3}\abs{a}+
        \dfrac{1}{4}\lambda\left(\lambda^2+2\abs{a}\lambda+2a^2+1\right),$$
	%$b_1=0$, $|b_2|<\dfrac{1}{2}$, 
    $$|b_3|\leq\dfrac{1}{3}\lambda^2+\dfrac{2}{3}\abs{a}\lambda$$
and $$|b_4|\leq\dfrac{1}{4}\lambda\left(\lambda^2+2\abs{a}\lambda+2a^2+1\right).$$
All of these inequalities are sharp.
\end{theorem}

\begin{proof} Suppose that $f=h+\bar{g}\in\mathcal{K}_{a,\lambda}$.
	It follows from \eqref{1} that
    \begin{equation}\label{T4-1}
        \left(1-b_1\right)z+\sum\limits_{n=2}^{\infty}(a_n-b_n)z^n=k_a(z)
    \end{equation}
    and
    \begin{equation}\label{T4-2}
        \sum\limits_{n=2}^{\infty}nb_nz^{n-1}=\sum\limits_{n=2}^{\infty}\lambda na_nz^n.
    \end{equation}
	By \eqref{T2-3}, \eqref{T3-3}, and \eqref{T3-4}, we get %$$k_a'(0)=1,$$ 
    $$\dfrac{k_a''(0)}{2!}=a,$$ 
    \begin{equation*}\label{T4-3}
        \dfrac{k_a'''(0)}{3!}=\dfrac{2}{3}a^2+\dfrac{1}{3}
    \end{equation*}
    and 
     \begin{equation*}\label{T4-4}
        \dfrac{{k_a}^{(4)}(0)}{4!}=\dfrac{1}{3}a^3+\dfrac{2}{3}a.
    \end{equation*}
	By comparing the coefficients of $z^2$, $z^3$ and $z^4$ on both sides of \eqref{T4-1}, we obtain
	\begin{equation}\label{T4-5}
		\begin{cases}
			%1-b_1=1,\\
			a_2-b_2=a,\\
			a_3-b_3=\frac{2}{3}a^2+\frac{1}{3},\\
			a_4-b_4=\frac{1}{3}a^3+\frac{2}{3}a.
		\end{cases}
	\end{equation}
	Similarly, it follows from \eqref{T4-2} that
	\begin{equation}\label{T4-6}
		\begin{cases}
			2b_2=\lambda,\\
			3b_3=2\lambda a_2,\\
			4b_4=3\lambda a_3.
		\end{cases}
	\end{equation}
   Combining \eqref{T4-5} and \eqref{T4-6} yields
    \begin{equation*}\label{T4-7}
        a_2=a+\dfrac{\lambda}{2},
    \end{equation*}
     \begin{equation*}\label{T4-8}
        a_3=\dfrac{1}{3}\left(\lambda^2+2a\lambda+2a^2+1\right),
    \end{equation*}
     \begin{equation*}\label{T4-9}
        a_4=\dfrac{1}{3}a^3+\dfrac{2}{3}a+
        \dfrac{1}{4}\lambda\left(\lambda^2+2a\lambda+2a^2+1\right),
    \end{equation*}
  
     \begin{equation*}\label{T4-10}
        b_3=\dfrac{1}{3}\lambda^2+\dfrac{2}{3}a\lambda
    \end{equation*}  and 
     \begin{equation*}\label{T4-11}
        b_4=\dfrac{1}{4}\lambda\left(\lambda^2+2a\lambda+2a^2+1\right).
    \end{equation*}
The inequalities in Theorem \ref{T4} follow directly from the above equations.
\end{proof}

\end{document}
